\section{数据工程 II：数据合成、难度课程与 RLVR 信号}

本讲中最值得注意的实践之一是数据合成（data synthesis）的地位提升。讲者明确表示，团队会把最强的人投入数据合成，而不是把它当成低价值流水线。这一判断与近年工业趋势一致。

在数学与代码任务中，合成数据不仅用于扩量，更用于 \textbf{补齐能力盲区}：当模型在若干子类任务表现弱，就按子类定向合成高质量样本并与原始数据混训，形成类似 curriculum learning 的提升路径。

\begin{importantbox}{课程中的关键现象：高数据效率}
讲者展示的经验是：在反馈可靠、任务定义清晰时，少量高质量示例也能触发显著学习增益。这并不否认大规模数据，而是强调 ``样本质量和任务覆盖'' 的边际价值更高。
\end{importantbox}

字幕中还提到 entropy collapse 问题：随着 RL 强化，策略可能过早收敛到窄分布，导致 pass@1 提升但 pass@k 不升反降。通过多样化合成和探索策略设计，可以在质量提升的同时维持候选解多样性。

\begin{table}[H]
\centering
\begin{tabular}{p{3.0cm}p{4.0cm}p{4.5cm}}
\toprule
\textbf{数据策略} & \textbf{直接收益} & \textbf{潜在风险} \\
\midrule
盲目扩量 & 覆盖面变宽 & 噪声样本稀释有效梯度 \\
难点定向合成 & 快速修补短板 & 过拟合合成模板 \\
原始+合成混训 & 稳定性更好 & 配比调参成本高 \\
仅追 pass@1 优化 & 指标短期提升快 & 探索能力衰减，pass@k 下降 \\
\bottomrule
\end{tabular}
\caption{数据合成策略与收益/风险平衡}
\end{table}

\begin{knowledgebox}{可验证合成的工程要点}
\begin{itemize}
    \item 题目必须可程序检验，避免伪标签漂移；
    \item 生成器与验证器分离，避免同源错误；
    \item 用对抗抽样抽查 ``看似正确但不可执行'' 的样本；
    \item 合成样本加入后应重新评估分布偏移。
\end{itemize}
\end{knowledgebox}

\subsection{本章小结}
数据合成不是简单扩充，而是能力结构设计。对 agent 训练而言，数据效率提升往往来自 ``更好的任务构造''，而不是 ``更多 token''。

